\section{Multi-homomorphisms}
\label{sec:multihom}

Let $\TT$ be an algebraic theory, and $\pair\CC\prod$ a category with
chosen finite products. For all objects $X,Y$ in $\CC$, and
natural number $k$, define the strength (for the reader
monad $X^k$):
\[
\str : X^k \times Y \xto{\seq[i=1][k]{\projection_i\times \id}}
  \parent{X\times Y}^k
\]
Let $\A_1, \ldots,
\A_n,\A[2] \in \Mod(\TT, \CC)$ be algebras. A function $f :
\prod_{i=1}^n\carrier{\A_i} \to \carrier{\A[2]}$ is a
\emph{multi-$\TT$-homomorphism}~\citep{hyland-power-discrete-lawvere-theories-journal}
when, for every term $t \in \TT_k$ and $1 \leq i_0 \leq n$, the
following diagram commutes: \diagram{01}
%
When $n = 1$, this condition degenerates to the following diagram,
which states that $f : \carrier\A \to \carrier{\A[2]}$ is a
$\TT$-homomorphism: \diagram{02}
%
Therefore, multi-$\TT$-homomorphisms generalise $\TT$-homomorphisms:
\begin{proposition}\propositionlabel{unary multi-TT-homomorphisms}
  Consider the bijection between $1$-ary multi-morphisms $h :
  (\carrier\A) \to \carrier\A$ and $\C$-morphisms $h' : \carrier\A
  \isomorphic \prod\set{\carrier\A} \xto{h} \carrier\A$.  A
  multi-morphism $h$ is a multi-$\TT$-homomorphism iff the
  corresponding $h'$ is a $\TT$-homomorphism.
\end{proposition}


% FIXME: maybe replace with the proof that multi-A-homomorphisms are multi-linear
% \begin{example}
%   Let $\Field$ be a field and $\Vect_{\Field}$ be the algebraic theory
%   of vector spaces over $\Field$. Consider the term $x_1, x_2 \vdash t
%   \definedby \lambda_1x_1 + \lambda_2x_2$ for some scalars $\lambda_1,
%   \lambda_2 \in \Field$. Then the multi-homomorphism condition, when
%   applied to vector spaces $\V[3], \V[1], \V[2]$, function $f :
%   \carrier{\V[3]} \times \carrier{\V[1]} \to \carrier{\V[2]}$ becomes
%   the usual bilinearity condition, i.e., for $i=1$ (and similarly for $i=2$):
%   \[
%   f(\lambda_1u_1 + \lambda_2u_2, v) = \lambda_1f(u_1, v) +
%   \lambda_2f(u_2,v)
%   \]
%   Since: \diagram{03} Therefore, the multi-linear maps, such as the
%   inner product $\pair-- : \Field^n\times\Field^n \to \Field$, are
%   multi-$\Vect_{\Field}$-homomorphisms.
% \end{example}
We will define a multi-category structure $\carrier- : \Multi(\TT,\C) \to
\pair\C\prod$ over the finite-product multi-category:
\begin{itemize}
\item The objects are $\TT$-algebras.
\item The multi-morphisms are the multi-$\TT$-homomorphisms.
\item The multi-functor $\carrier-$ sends an algebra to its carrier
  and a multi-$\TT$-homomorphism to its underlying morphism.
\item The identities are the identity morphisms.
\item Composition is inherited from $\pair\C\prod$.
\end{itemize}
The next few propositions show that this data defines a
multi-category.

\begin{proposition}\propositionlabel{multi-TT-homomorphism composition}
  Let $g : (\A[2]_1, \ldots, \A[2]_n) \to \A[3]$ and $f_i :
  (\A[1]_{i1}, \ldots, \A[1]_{im_i}) \to \A[2]_i$ be
  multi-$\TT$-homomorphisms. Then $g \compose (f_1, \ldots, f_n) :
  (\A[1]_{11}, \ldots, \A[1]_{nm_n}) \to \A[3]$ is a
  multi-$\TT$-homomorphism.
\end{proposition}
\begin{proof}
  Let $p := \sum_{i=1}^nm_i$ and $\phi : [p] \xto{\isomorphic}
  \coprod_{i=1}^n[m_i]$ be the enumeration from the definition of
  composition in $\pair\C\prod$ for this case. Take any $1 \leq \ell_0
  \leq p$, and let $\pair{i_0}{j_0} := \phi\ell_0$ be its
  corresponding pair of indices. Take any $k$-ary term $t \in \TT_k$,
  and calculate as in \figref{multi-hom composition}.
  \begin{figure}
    \diagram{04}
    \caption{Multi-$\TT$-homomorphism property of the composition}\figlabel{multi-hom composition}
  \end{figure}
\end{proof}

\begin{proposition}\propositionlabel{Multi(TT,C) multi-category}
  The data above defines a multi-category $\Multi(\TT, \C)$ together
  with a faithful multi-functor $\carrier- : \Multi(\TT,\C) \to
  \pair\C\prod$: $\carrier{f_1} = \carrier{f_2} \implies f_1 = f_2$
  for every parallel pair of multi-$\TT$-homomorphisms $f_1,f_2 :
  (\A_1, \ldots, \A_n) \to \A[2]$.
\end{proposition}
\begin{proof}
  Identities are multi-$\TT$-homomorphisms by \propositionref{unary
    multi-TT-homomorphisms}, and composition of
  multi-$\TT$-homomorphisms is a multi-$\TT$-homomorphism by
  \propositionref{multi-TT-homomorphism composition}, and so we have a
  multi-category structure.

  Since composition and identities are inherited from $\pair\C\prod$,
  we have the structure of multi-functor. This multi-functor structure
  is faithful since the multi-morphisms in $\Multi(\TT, \C)$ are those
  of $\pair\C\prod$ with additional structure. Since faithful
  multi-functors reflect multi-categories, we conclude.
\end{proof}
% Let $\ppCAT$ be the $2$-category consisting of:
% \begin{itemize}
% \item $0$-cells are $\pair\C\prod$, categories with chosen finite products;
% \item $1$-cells are finite product preserving functors, not
%   necessarily preserving the chosen products;
% \item $2$-cells are the natural transformations between these functors.
% \end{itemize}
% This $2$-category is a sub-category of $\CAT$ that is full on
% $2$-cells but not on $1$-cells. This seems a bit weird, but that's the
% structure we seem to need. Given such a product preserving functor $H
% : \pair \C\prod \to \pair{\C[2]}\prod$, by fiat we have canonical
% isomorphisms $H\prod_{i=1}^nA \isomorphic \prod_{i=1}^nHA$ that
% preserve the product structure.

% We define the data of a $2$-functor $\Multi(\TT, -) : \ppCAT \to
% \MultiCAT$ sending:
% \begin{itemize}
% \item Each category with chosen products ($0$-cell) $\C$ to the
%   multi-category ($0$-cell) $\Multi(\TT, \C)$.
% \item Each finite-product preserving functor ($1$-cell) $H : \C[1] \to
%   \C[2]$ to the following multi-functor $\Multi(\TT, H) : \Multi(\TT,
%   \C[1]) \to \Multi(\TT, \C[2])$, sending:
%   \begin{itemize}
%   \item each $\TT$-algebra $\A$ to the $\TT$-algebra $\Multi(\TT,
%     H):=H\A$, given by:
%     \begin{itemize}
%     \item its carrier $\carrier{H\A}$ is $H\carrier{\A}$; and
%     \item the interpretation of each $t \in \TT_n$ is $H\A\sem t :
%       (H\carrier\A)^n \xto\isomorphic H(\carrier\A^n) \xto{\A\sem t}
%       H\carrier\A$.
%     \end{itemize}
%     More abstractly, an algebra is a product-preserving functor $\A :
%     \TT \to \C$, and $H\A$ is the product preserving functor $H\A :
%     \TT\xto\A \C \xto H \C[2]$.
%   \item each multi-$\TT$-homomorphism $h : (\A_1, \ldots, \A_n) \to
%     \A[2]$ to the multi-$\TT$-homomorphism $\Multi(\TT,H)h : (H\A_1,
%     \ldots, H\A_n) \to H\A_n$, given by:
%     \[
%     \Multi(\TT,H)h : \prod_{i=1}^nH\A_i \xto{\isomorphic}
%     H(\prod_{i=1}^n\A_i) \xto{Hh} H\A[2]
%     \]
%   \end{itemize}
% \item Each natural transformation $\quiver{nat trans between pp
%   functors}$ ($2$-cell) to the multi-natural transformation
%   $\quiver{Mod(TT, alpha) type}$, given, at an algebra $\A\in
%   \Mod(\TT,\C[1])$, by the unary multi-homomorphism corresponding to
%   the following $\TT$-homomorphism (cf.~\propositionref{unary
%     multi-TT-homomorphisms}):
%   \[
%   \Multi(\TT,\alpha)_{\A}' : H\carrier{\A} \xto{\alpha_{\carrier \A}} K\carrier{\A}
%   \]
% \end{itemize}
% It may be worthwhile to spell-out some additional related
% structure. The following forms the data of a $2$-functor $\seq- :
% \ppCAT \to \MultiCAT$, sending:
% \begin{itemize}
% \item $0$-cells: categories with specified products $\pair\C\prod$ to
%   their associated multi-category $\pair\C\prod$
%   (cf.~\exampleref{multi-cat from chosen products}).
% \item $1$-cells: finite-product preserving functors $H : \C \to
%   \C[2]$ to the multi-functor sending:
%   \begin{itemize}
%   \item objects $a \in \Obj{\pair\C\prod} = \Obj\C$ to $\seq Ha \in \Obj{\C[2]}$;
%   \item multi-morphisms $f:(a_1, \ldots, a_n) \to b$ to the multimorphism:
%     \[
%     \seq Hf : \prod_{i=1}^n Ha_i \xto{\isomorphic} H(\prod_{i=1}^na_i) \xto{Hf} Hb
%     \]
%   \end{itemize}
% \item $2$-cells: natural transformations $\quiver{nat trans between pp
%   functors}$ to the multi-natural transformation given, for each $a
%   \in \Obj\C$, by:
%   \[
%   \seq[a]{\alpha}{} : \prod_{i=1}^1Ha \isomorphic Ha \xto{\alpha_a} Ka
%   \]
% \end{itemize}
% This data allows us to talk about the following data for a natural
% $2$-transformaion $\quiver{forgetful functor 2-transformation: type}$
% whose components, for each category $\pair\C\prod$ with chosen
% finite-products, consists of the forgetful multi-functors:
% \begin{itemize}
% \item sending each $\TT$-algebra $\A$ to its carrier $\carrier\A$;
% \item sending each multi-$\TT$-homomorphism $h : (\A_1, \ldots, \A_n)
%   \to \A[2]$ to its underlying morphism $\carrier h : \prod_{i=1}^n
%   \carrier\A_i \to \carrier{\A[2]}$.
% \end{itemize}
% \begin{proposition}
%   This data defines:
%   \begin{itemize}
%   \item Two $2$-functors $\seq-, \Multi(\TT,-) : \ppCAT \to
%     \MultiCAT$, the former, $\seq-$, is $1$-faithful and
%     $2$-fully-faithful.
%   \item A natural $2$-transformation $\quiver{forgetful functor
%     2-transformation: type}$ whose components are faithful
%     multi-functors.
%   \end{itemize}
% \end{proposition}
% \begin{proof}
%   We have much to do. Starting with the well-formedness of the data
%   for $\seq-$, we have shown $\pair\C\prod$ is a multi-category. To
%   see that $\seq H$ is a functor, calculate as in \figref{seq h functoriality}.
%   \begin{figure}
%   \diagram{05}
%   \diagram{06}
%   \caption{Functoriality of $\seq H$}\figlabel{seq h functoriality}
%   \end{figure}

%   Next, we'll use the following observation. Given a natural
%   transformation $\quiver{nat trans between pp functors}$ between
%   finite-product preserving functors:
%   \[
%     \diagram*{07}\tag*{($*$)}
%   \]
%   proved by inverting the canonical isomorphisms and using the
%   universal property of products:
%   \diagram{08}
%   Using this property, we show that $\quiver{seq alpha type}$ is a
%   multi-natural transformation. Take any multi-morphism $f : (a_1,
%   \ldots, a_n) \to b$, and calculate:
%   \diagram{09}
%   Therefore the data for the $2$-functor $\seq-$ is well-defined.

%   Next, consider the data for the $2$-functor $\Mod(\TT,-)$.
%   \propositionref{Multi(TT,C) multi-category} shows it is a
%   multi-category, and so we need to check the well-formedness of the
%   higher dimensional data. For $1$-cells, as we explained more
%   abstractly, the data $\Multi(\TT, H)\A$ forms an algebra.  Given a
%   multi-$\TT$-homomorphism $h : (\A_1, \ldots, \A_n) \to \A[2]$, we
%   need to show that the morphism:
%   \[
%   \Multi(\TT, H)h : (H\A_1, \ldots, H\A_n) \to H\A[2]
%   \]
%   is indeed a multi-$\TT$-homomorphism. Take any $t \in
%   \TT_k$, $i_0 < n$, and calculate as in \figref{multi-TT-homomorphism
%     of Multi(TT, H)h}.

%   \begin{figure}
%     \diagram{10}
%     \caption{Multi-homomorphism proof for $\Multi(\TT,H)h$
%     }\figlabel{multi-TT-homomorphism of Multi(TT, H)h}
%   \end{figure}

%   To show that the structure on $2$-cells is well-defined, we need to
%   show that $\Mod(\TT,\alpha)_{\A}' : H\A \to K\A$ is a
%   $\TT$-homomorphism. Take any $t \in \TT_k$, and calculate:
%   \diagram{11}

%   Checking that $\seq-$ is a $2$-functor is straightforward:
%   \begin{itemize}
%   \item For $1$-cells, take two finite-product preserving functors
%     $\C[3] \xto{H} \C[4] \xto{K} \C[1]$.
%     \begin{itemize}
%     \item For objects, we have $\seq{K \compose H}a := KH a := \seq K
%       \compose \seq H a$.
%     \item For multi-morphisms $f : (a_1, \ldots, a_n) \to b$,
%       calculate:
%       \diagram{12}
%     \end{itemize}
%     It's also straightforward to check that $\seq\Id$ is the identity
%     multi-functor:
%     \diagram{15}
%   \item For $2$-cells:
%     \begin{itemize}
%     \item For the identity natural transformation we have:
%     \[
%     \seq[a]{\id[H]} : \prod\set{Ha} \xto\isomorphic Fa = Fa
%     \]
%     \item For vertical composition, calculate for vertically
%       composable $2$-cells $\quiver{pp vertical composition}$:
%       \diagram{13}
%     \item For horizontal composition, calculate for $\quiver{pp
%       horizontal composition}$:
%     \diagram{14}
%     \end{itemize}
%   \end{itemize}
%   And thus we have a $2$-functor $\seq- : \ppCAT \to \MultiCAT$.

%   Finally, we turn to the natural $2$-transformation $\carrier - :
%   \Multi(\TT,-) \to \seq{-}$.
%   \begin{itemize}
%   \item For the $1$-naturality condition, take a product preserving
%     functor $H : \C[1] \to \C[2]$ in $\ppCAT$ and show the
%     following equality of multi-functors:
%     \diagram{16}
%     by chasing an algebra $\A \in \Multi(\TT, \C[1])$ (top) and a
%     multi-$\TT$-homomorphism $h : (\A_1, \ldots, \A_n)\to \A[2]$
%     (bottom):
%     \begin{center}
%       \diagram*{17}
%       \hfill
%       \diagram*{18}
%     \end{center}
%   \item For the $2$-naturality condition, take a natural
%     transformation $\quiver{nat trans between pp functors}$ and
%     algebra $\A$, and consider the component at $\A$ of the
%     appropriate pastings (drawn badly :D):
%     \[
%     \quiver{carrier 2-naturality}
%     \]
%   \end{itemize}
%   And so we have a natural $2$-transformation $\carrier- :
%   \Multi(\TT,-) \to \seq-$ \emph{provided} we show that $\Multi(\TT,-)
%   : \ppCAT \to \MultiCAT$ is indeed a $2$-functor. We can deduce that
%   from the fact that $\carrier-$ is component-wise injective:
%   \begin{itemize}
%   \item For identity $1$-cells:
%     \begin{itemize}
%     \item On objects, we indeed check that the carrier and algebra
%       structure of $\Multi(\TT, \Id)\A$ are the same as of $\A$.
%     \item On multi-homormophisms, we use the faithfulness of
%       $\carrier-_{\C}$:
%       \[
%       \carrier-_{\C}\compose \Multi(\TT, \Id) h
%       \explain={$1$-cell $2$-naturality}
%       \seq\Id \compose \carrier-_{\C}h
%       \explain*={$\seq-$ is a $2$-functor}
%       \carrier-_{\C}h
%       =
%       \carrier-_{\C}\compose \Id h
%       \]
%       and so by faithfulness, $\Multi(\TT, \Id)h = \Id h$.
%     \end{itemize}
%   \item Similarly, for composable $1$-cells $\C[3] \xto H \C[4] \xto K \C$:
%     \begin{itemize}
%     \item On objects, we indeed check that the algebras are the same.
%     \item On multi-homomorphisms, we use faithfulness:
%       \[
%       \carrier{\Multi(\TT, K \compose H)h}_{\C}
%       \explain={$1$-cell $2$-naturality}
%       \seq{K \compose H} (\carrier{h}_{\C[3]})
%       \explain*={$\seq-$ is a $2$-functor}
%       \seq{K} \compose \seq{H} (\carrier{h}_{\C[3]})
%       =
%       \carrier{\Multi(\TT, K)\compose \Multi(\TT, H)h}_{\C}
%       \]
%     \end{itemize}
%     and so $\Mod(\TT,-)$ is functorial on $1$-cells.
%   \item For $2$-cell identities, the same technique applies:
%     \[
%     \carrier{\Multi(\TT, \id[H])_{\A}}_{\C[2]}
%     \explain={$2$-cell $2$-naturality}
%     \seq[{\carrier\A}]{\id[H]}
%     \explain*={$\seq-$ is a $2$-functor}
%     (\id[\seq{H}])_{\carrier\A}
%     =
%     (\id[\seq{H}{\carrier\A}])
%     \explain={$1$-cell $2$-naturality}
%     \id[\carrier{\Multi(\TT, H)\A}]
%     \explain*={$\carrier-_{\C[2]}$ multi-functoriality}
%     \carrier{\id[\Multi(\TT, H)_{\A}]}
%     \]
%   \item For vertical composition, take a composable pair of $2$-cells,
%     and calulate as follows (reverting to the more succinct notation
%     in the definition of natural $2$-transformations):
%     \begin{gather*}
%       \quiver{mult 2-functoriality vertical composition: a}\\
%     \begin{aligned}
%       &\explain={identities}
%       \quiver{mult 2-functoriality vertical composition: b}
%       \explain={$2$-cell\\$2$-naturality}
%       \quiver{mult 2-functoriality vertical composition: c}
%       \explain={$G := \seq-$ is a\\$2$-functor}
%       \quiver{mult 2-functoriality vertical composition: d}\\
%       &\explain={$2$-cell\\$2$-naturality}
%       \quiver{mult 2-functoriality vertical composition: e}
%       \explain={$2$-cell\\$2$-naturality}
%       \quiver{mult 2-functoriality vertical composition: f}
%       \explain={identities}
%       \quiver{mult 2-functoriality vertical composition: g}
%     \end{aligned}
%     \end{gather*}
%   \item For horizontal composition, we use the same technique:
%     \begin{gather*}
%       \quiver{mult 2-functoriality horizontal composition: a}\\
%       \begin{aligned}
%         &\explain={identities}
%         \quiver{mult 2-functoriality horizontal composition: b}
%         \explain={$2$-cell\\$2$-naturality}
%         \quiver{mult 2-functoriality horizontal composition: c}\\
%         &\explain={$G$ is a\\$2$-functor}
%         \quiver{mult 2-functoriality horizontal composition: d}
%         \explain={identities}
%         \quiver{mult 2-functoriality horizontal composition: e}\\
%         &
%         \explain={$2$-naturality}
%         \quiver{mult 2-functoriality horizontal composition: f}
%         \explain={$2$-naturality}
%         \quiver{mult 2-functoriality horizontal composition: g}\\
%         &\explain={identities}
%         \quiver{mult 2-functoriality horizontal composition: h}
%       \end{aligned}
%     \end{gather*}
%   \end{itemize}
%   Therefore, $\Mod(\TT,-)$ is also a $2$-functor.
% \end{proof}
